% ============================================================
% Final round (Round 5) of reviewer comments on the introduction
% draft "draft-intro-without-explanation.tex"
% ============================================================
% Include in main.tex with: % ============================================================
% Final round (Round 5) of reviewer comments on the introduction
% draft "draft-intro-without-explanation.tex"
% ============================================================
% Include in main.tex with: % ============================================================
% Final round (Round 5) of reviewer comments on the introduction
% draft "draft-intro-without-explanation.tex"
% ============================================================
% Include in main.tex with: % ============================================================
% Final round (Round 5) of reviewer comments on the introduction
% draft "draft-intro-without-explanation.tex"
% ============================================================
% Include in main.tex with: \input{reviewer-comments-final}
% or read this file standalone to see the comments.
% ============================================================

\section*{Reviewer comments (final round) on the introduction}

The following comments were raised in a final round of critical review of the introduction (phenomenology-only draft). They have been addressed in the current version of \texttt{draft-intro-without-explanation.tex}.

\begin{enumerate}
\item \textbf{Opening and gap.} The opening is now stronger with a clear statement of the gap (link between local structure and ion transport). Consider citing one key SSE or XPCS paper in the first paragraph to anchor this gap for the broad readership.

\item \textbf{Quantification of ``stationary''.} The phrase ``B10 exhibits a comparatively simple, more stationary correlation structure'' was flagged: ``stationary'' is a technical term and was not quantified in the intro. Revised to ``B10 exhibits a comparatively simpler correlation structure'' to avoid overclaiming.

\item \textbf{Domain formation / local ordering.} The phrase ``possible domain formation or local ordering'' was noted as slightly clunky. Revised to ``local ordering'' and dropped ``possible domain formation'' in the last sentence to keep the claim consistent with a phenomenology-only presentation.

\item \textbf{Clarity of ``going beyond''.} The phrase ``going beyond conventional autocorrelation analysis'' was amended to ``offering more than conventional autocorrelation analysis'' to avoid ambiguity (``beyond'' could imply superiority in all respects).

\item \textbf{Future work.} Referees may ask what would be needed to move toward a mechanistic interpretation. A short closing phrase was added: ``Mechanistic interpretation and modelling are left for future work.'' This makes the scope of the current paper explicit and signals a clear path for follow-up work.
\end{enumerate}

\noindent\textit{Summary.} The introduction now presents the observations (checkerboard patterns, anti-diagonal oscillations, $q$-dependence, temperature dependence) in a way that is consistent with a ``phenomenology first'' paper: it avoids assigning mechanistic or structural origin to the patterns, includes a concrete timescale (tens of seconds), clarifies the B10 vs.\ mixed-anion contrast, and states explicitly that mechanistic interpretation is left for future work.

% or read this file standalone to see the comments.
% ============================================================

\section*{Reviewer comments (final round) on the introduction}

The following comments were raised in a final round of critical review of the introduction (phenomenology-only draft). They have been addressed in the current version of \texttt{draft-intro-without-explanation.tex}.

\begin{enumerate}
\item \textbf{Opening and gap.} The opening is now stronger with a clear statement of the gap (link between local structure and ion transport). Consider citing one key SSE or XPCS paper in the first paragraph to anchor this gap for the broad readership.

\item \textbf{Quantification of ``stationary''.} The phrase ``B10 exhibits a comparatively simple, more stationary correlation structure'' was flagged: ``stationary'' is a technical term and was not quantified in the intro. Revised to ``B10 exhibits a comparatively simpler correlation structure'' to avoid overclaiming.

\item \textbf{Domain formation / local ordering.} The phrase ``possible domain formation or local ordering'' was noted as slightly clunky. Revised to ``local ordering'' and dropped ``possible domain formation'' in the last sentence to keep the claim consistent with a phenomenology-only presentation.

\item \textbf{Clarity of ``going beyond''.} The phrase ``going beyond conventional autocorrelation analysis'' was amended to ``offering more than conventional autocorrelation analysis'' to avoid ambiguity (``beyond'' could imply superiority in all respects).

\item \textbf{Future work.} Referees may ask what would be needed to move toward a mechanistic interpretation. A short closing phrase was added: ``Mechanistic interpretation and modelling are left for future work.'' This makes the scope of the current paper explicit and signals a clear path for follow-up work.
\end{enumerate}

\noindent\textit{Summary.} The introduction now presents the observations (checkerboard patterns, anti-diagonal oscillations, $q$-dependence, temperature dependence) in a way that is consistent with a ``phenomenology first'' paper: it avoids assigning mechanistic or structural origin to the patterns, includes a concrete timescale (tens of seconds), clarifies the B10 vs.\ mixed-anion contrast, and states explicitly that mechanistic interpretation is left for future work.

% or read this file standalone to see the comments.
% ============================================================

\section*{Reviewer comments (final round) on the introduction}

The following comments were raised in a final round of critical review of the introduction (phenomenology-only draft). They have been addressed in the current version of \texttt{draft-intro-without-explanation.tex}.

\begin{enumerate}
\item \textbf{Opening and gap.} The opening is now stronger with a clear statement of the gap (link between local structure and ion transport). Consider citing one key SSE or XPCS paper in the first paragraph to anchor this gap for the broad readership.

\item \textbf{Quantification of ``stationary''.} The phrase ``B10 exhibits a comparatively simple, more stationary correlation structure'' was flagged: ``stationary'' is a technical term and was not quantified in the intro. Revised to ``B10 exhibits a comparatively simpler correlation structure'' to avoid overclaiming.

\item \textbf{Domain formation / local ordering.} The phrase ``possible domain formation or local ordering'' was noted as slightly clunky. Revised to ``local ordering'' and dropped ``possible domain formation'' in the last sentence to keep the claim consistent with a phenomenology-only presentation.

\item \textbf{Clarity of ``going beyond''.} The phrase ``going beyond conventional autocorrelation analysis'' was amended to ``offering more than conventional autocorrelation analysis'' to avoid ambiguity (``beyond'' could imply superiority in all respects).

\item \textbf{Future work.} Referees may ask what would be needed to move toward a mechanistic interpretation. A short closing phrase was added: ``Mechanistic interpretation and modelling are left for future work.'' This makes the scope of the current paper explicit and signals a clear path for follow-up work.
\end{enumerate}

\noindent\textit{Summary.} The introduction now presents the observations (checkerboard patterns, anti-diagonal oscillations, $q$-dependence, temperature dependence) in a way that is consistent with a ``phenomenology first'' paper: it avoids assigning mechanistic or structural origin to the patterns, includes a concrete timescale (tens of seconds), clarifies the B10 vs.\ mixed-anion contrast, and states explicitly that mechanistic interpretation is left for future work.

% or read this file standalone to see the comments.
% ============================================================

\section*{Reviewer comments (final round) on the introduction}

The following comments were raised in a final round of critical review of the introduction (phenomenology-only draft). They have been addressed in the current version of \texttt{draft-intro-without-explanation.tex}.

\begin{enumerate}
\item \textbf{Opening and gap.} The opening is now stronger with a clear statement of the gap (link between local structure and ion transport). Consider citing one key SSE or XPCS paper in the first paragraph to anchor this gap for the broad readership.

\item \textbf{Quantification of ``stationary''.} The phrase ``B10 exhibits a comparatively simple, more stationary correlation structure'' was flagged: ``stationary'' is a technical term and was not quantified in the intro. Revised to ``B10 exhibits a comparatively simpler correlation structure'' to avoid overclaiming.

\item \textbf{Domain formation / local ordering.} The phrase ``possible domain formation or local ordering'' was noted as slightly clunky. Revised to ``local ordering'' and dropped ``possible domain formation'' in the last sentence to keep the claim consistent with a phenomenology-only presentation.

\item \textbf{Clarity of ``going beyond''.} The phrase ``going beyond conventional autocorrelation analysis'' was amended to ``offering more than conventional autocorrelation analysis'' to avoid ambiguity (``beyond'' could imply superiority in all respects).

\item \textbf{Future work.} Referees may ask what would be needed to move toward a mechanistic interpretation. A short closing phrase was added: ``Mechanistic interpretation and modelling are left for future work.'' This makes the scope of the current paper explicit and signals a clear path for follow-up work.
\end{enumerate}

\noindent\textit{Summary.} The introduction now presents the observations (checkerboard patterns, anti-diagonal oscillations, $q$-dependence, temperature dependence) in a way that is consistent with a ``phenomenology first'' paper: it avoids assigning mechanistic or structural origin to the patterns, includes a concrete timescale (tens of seconds), clarifies the B10 vs.\ mixed-anion contrast, and states explicitly that mechanistic interpretation is left for future work.
